\documentclass[10pt]{article}
% Shared look for overview.tex and project-details.tex. Build with: tectonic -X compile <file>.tex
\usepackage[a4paper,margin=1.7cm]{geometry}
\usepackage{fontspec}
\setmainfont{texgyrepagella}[Extension=.otf,UprightFont=*-regular,BoldFont=*-bold,ItalicFont=*-italic,BoldItalicFont=*-bolditalic]
\setsansfont{texgyreheros}[Extension=.otf,UprightFont=*-regular,BoldFont=*-bold,ItalicFont=*-italic,BoldItalicFont=*-bolditalic]
\usepackage{amsmath,microtype,xcolor,tikz,booktabs,tabularx,enumitem,titlesec,array}
\usepackage[hidelinks]{hyperref}
\usetikzlibrary{positioning,arrows.meta,fit}

\definecolor{ink}{HTML}{1F2933}
\definecolor{accent}{HTML}{1F6F8B}
\definecolor{soft}{HTML}{E8F1F4}
\definecolor{muted}{HTML}{5C6B73}
\color{ink}

\titleformat{\section}{\sffamily\bfseries\large\color{accent}}{}{0pt}{}
\titlespacing*{\section}{0pt}{1.1ex}{0.5ex}
\titleformat{\subsection}{\sffamily\bfseries\normalsize}{}{0pt}{}
\titlespacing*{\subsection}{0pt}{1ex}{0.3ex}
\setlist{itemsep=1pt,topsep=2pt,parsep=0pt,leftmargin=1.2em}
\setlength{\parindent}{0pt}
\setlength{\parskip}{0.5ex}
\pagestyle{empty}

\newcommand{\code}[1]{\texttt{\small #1}}
% \gh[branch]{path}{text}: a link to that place in the group repo on GitHub (default branch main)
\newcommand{\repo}{https://github.com/angelabuxo/IMDb-movie-success}
\newcommand{\gh}[3][main]{\href{\repo/tree/#1/#2}{\color{accent}#3}}
\newcommand{\where}[1]{{\small\color{muted}#1}}
\newcolumntype{L}{>{\raggedright\arraybackslash}X}

\newcommand{\doctitle}[2]{%
  {\sffamily\bfseries\LARGE\color{accent}#1\par}\vspace{2pt}
  {\small\color{muted}#2\par}\vspace{4pt}}


\begin{document}
\doctitle{Project details}{Log-Log Legends · DTU 02805 · state on 14 November 2026 · the one-page version is \textit{overview.pdf}. Each analysis gets a short block: what it asks, where it lives, what it found.}

\section{Data and shared rules}

\begin{tabularx}{\linewidth}{@{}p{5.4cm}rL@{}}
\toprule
What & Count & Source \\
\midrule
Movies with 10,000+ votes (the graph) & 12,673 & IMDb \\
Clean films with budget and gross in USD & 2,446 (100k+), 6,851 (10k+) & Wikipedia infoboxes \\
Screenplays & 1,095 clean 100k+ films & MovieSum \\
Audience reviews & 2,070 clean 100k+ films, to 2017 & MovieLens Tag Genome \\
Critic scores & Metacritic 2,431, RT 2,740 & Wikipedia, Wikidata \\
\bottomrule
\end{tabularx}

\smallskip
``Clean'' means budget and gross both in USD and at least \$1,000, no rentals-only or regional gross, no streaming
distributor, and released before 2026. The 100k+ tier favours hits: only 169 of its 2,446 films grossed less than their
budget. Profit is $\ln(\text{gross}/\text{budget})$ everywhere. Network measures use only earlier films, every result
gets a null or a second method, and no script text, review text or IMDb fields go into git.

\section{Who made the film}

\subsection{Movie graph \where{(\code{scripts/network/}, \code{notebooks/movie\_graph.ipynb})}}
Movies, people (by role and billing) and companies in one network, with studio names merged through Wikidata IDs.
\begin{itemize}
  \item The studio's track record is the strongest single measure: $\rho = +0.27$ with the multiple, net of budget and year, on 4,568 films.
  \item All network measures together lift $R^2$ from 0.041 to 0.130 over budget, year and genre (shuffle $p = 0.005$).
  \item A director's next film is a hit 35\% of the time after a hit, and 18\% after a non-hit.
\end{itemize}

\subsection{Shared-people returns \where{(\code{scripts/network/roi\_similarity.py})}}
Do films with the same director, producer, cast member or studio have more similar returns than chance?
\begin{itemize}
  \item Every role clusters: 0.09 (top-three cast) to 0.17 (producers), against a shuffle null of at most 0.02.
  \item It survives dropping franchise films, ranking the returns, and keeping only films from 1980 on. About half survives a shuffle within each studio.
  \item It's descriptive: it uses all years at once.
\end{itemize}

\subsection{Actor network \where{(Niklas, branch \code{initial\_actor\_network}, not merged)}}
Top-three billed actors, linked by shared films. It comes with an interactive page with projections,
communities, centrality and per-actor ROI.

\section{What the film says}

\subsection{Scripts and reviews notebook \where{(\code{notebooks/scripts\_and\_reviews.ipynb})}}
\begin{itemize}
  \item Only 2 of 15 script features pass $q < 0.05$ net of budget and year. The strongest is a character network centred on one character ($\rho = -0.097$).
  \item Six VADER arc shapes don't differ in profit ($p = 0.99$).
\end{itemize}

\subsection{Scattertext words \where{(\code{analysis/scripts\_text/}, pre-registered)}}
Top against bottom quarter by multiple, within genre and decade, on dialogue.
\begin{itemize}
  \item The difference is slightly above random groupings ($z = 2.08$, $p = 0.024$).
  \item It disappears without franchise films ($z = 0.50$). Dropping the same number of random films never does that (0 of 50 tries), so the signal sits in the franchise films themselves.
  \item A word classifier doesn't beat genre, decade and budget (AUC 0.58 against 0.58).
\end{itemize}

\subsection{Emotional arcs \where{(\code{analysis/arcs/}, pre-registered, one test)}}
Every sentence of 834 non-franchise scripts is scored by DistilBERT SST-2. Each film becomes a 100-point tone
curve, summarised as swing (D1) and ending lift (D2). Budget, genre, decade and scene length are controlled.
\begin{itemize}
  \item Swing: no link ($\rho = +0.035$, $p = 0.32$).
  \item Ending lift: a near miss ($\rho = +0.078$, $p = 0.028$, against a pre-set cutoff of 0.025). Its sign holds in 12 of 13 variants, but not with VADER.
  \item The arc adds nothing to out-of-sample prediction.
  \item About half of the arc variation between films is scene noise (split-half reliability around 0.45).
  \item A fourth sentiment model, ModernBERT, is running as a robustness check.
\end{itemize}

\section{How it was received}

\subsection{Reviews exploration \where{(Àngela, branch \code{reviews-exploration}, not merged)}}
2,095 films up to 2017. Every correlation is tested against a shuffle within genre $\times$ decade.
\begin{itemize}
  \item The multiple goes with IMDb rating at $\rho = 0.27$, Metacritic 0.26 and Rotten Tomatoes 0.33. Genre and decade alone give about 0.1.
  \item The critics agree with each other far more (Metacritic and Rotten Tomatoes at 0.93). Review tone goes with the multiple at 0.13.
  \item A disagreement score (rating rank minus money rank) finds ``liked beyond their money'' films (\textit{tarantino}, \textit{noir}) and ``earned beyond their liking'' films (\textit{scares}, \textit{slasher}).
\end{itemize}

\section{What predicts profit \where{(\code{notebooks/what\_predicts\_profit.ipynb})}}
Trained on films to 2010, tested on 2011 to 2017.

\begin{tabularx}{0.62\linewidth}{@{}Lr@{}}
\toprule
Model & Out-of-sample $R^2$ \\
\midrule
Budget, year, genre & $-0.010$ \\
+ network & 0.202 \\
All pre-release blocks (adds month, franchise) & 0.244 \\
+ critic and audience scores & 0.312 \\
\bottomrule
\end{tabularx}
\begin{itemize}
  \item On films with a screenplay, script features add nothing beyond the network ($R^2$ 0.037 against 0.036).
  \item Director hit streaks beat a shuffle of each director's own films (+0.175 against +0.090). PageRank adds nothing beyond plain experience.
  \item Rotten Tomatoes goes with the multiple at $\rho = +0.25$. There are 158 loved flops and 205 panned hits.
\end{itemize}

\section{Open and weak spots}
\begin{itemize}
  \item \textbf{Human validation}, which only people can fill:
  \begin{itemize}
    \item Scattertext: \code{analysis/scripts\_text/outputs/validation\_*.csv}
    \item Arcs: \code{analysis/arcs/validation/} and \code{analysis/arcs/outputs/s12\_back\_to\_text.csv}
    \item Reviews: the 50-review sample on \code{reviews-exploration}
  \end{itemize}
  \item \textbf{Arcs Step 11} waits on the instrument question (whole-review VADER against sentence-level VADER).
  \item \textbf{Weak spots}: the franchise flag is rough, release month isn't checked, reviews stop at 2017, and gross over budget isn't profit.
\end{itemize}

\section{Where things are}
\begin{tabularx}{\linewidth}{@{}p{5.6cm}L@{}}
\toprule
\code{scripts/} & Data download and parsing, the graph, network features \\
\code{analysis/scripts\_text/}, \code{analysis/arcs/} & The two pre-registered text analyses: config, LOG, steps, outputs \\
\code{notebooks/} & Graph, profit model, scripts and reviews, one folder per text analysis \\
\code{site/scripts/}, \code{site/arcs/} & Static pages for the text analyses \\
\code{data/bundles/} & Compressed data that ships with the repo \\
\bottomrule
\end{tabularx}
\end{document}
